\Needspace{18\baselineskip}
\subsection{A countable divisorial positivity criterion}
\label{subsec:countable-divisorial-positivity}

A prescribed divisorial part can determine an ideal envelope's
ordinary type even when it has infinitely many components.
On a pluriclosed projective surface, its smooth residual then
determines the mass of every representative of that type.
This gives a positivity test directly from divisorial orders.

Let $X=\mathbb P^2$, and let $\kappa$ be the Fubini--Study form
with $\int_X\kappa^2=1$. Choose distinct hyperplanes $H_j$ and
coefficients $a_j>0$ with $A=\sum_{j\geq1}a_j<\infty$.
For unit defining sections $s_j$, put
\[
 \ell_j(Z)=\log\frac{|s_j(Z)|^2}{\|Z\|^2}\leq0,
 \qquad u_0=\sum_{j\geq1}a_j\ell_j,
 \qquad u=u_0-\sup_Xu_0.
\]
Let $\theta$ be a smooth real $(1,1)$-form satisfying
$\ddc\theta=0$ and $R=\theta-A\kappa>0$.

\Needspace{14\baselineskip}
\begin{lemma}\label{countable-divisorial:profile}\index{Divisorial current!summable hyperplane series}
The series defines a quasi-plurisubharmonic function, and
\[
 \ddc u=S-A\kappa,\qquad
 S=\sum_{j\geq1}a_j[H_j],\qquad
 \theta+\ddc u=R+S>0.
\]
For every prime divisor $E$ on every smooth compact modification
$\pi:Z\to X$, the convergent order profile is
\[
 b(E):=\sum_{j\geq1}a_j\operatorname{ord}_E(\pi^*H_j)
       =\nu(u\circ\pi,E).
\]
\end{lemma}
\begin{proof}
The Poincar\'e--Lelong formula gives
$\ddc\ell_j=[H_j]-\kappa$. Unitary invariance reduces its
integral to the distribution of $x=|Z_0|^2/\|Z\|^2$, whose
density on $[0,1]$ is $2(1-x)$. Thus
\[
 \int_X\ell_j\kappa^2=\int_0^1 2(1-x)\log x\,\mathrm dx=-\frac32.
\]
The decreasing series has finite $L^1$ norm and a common
lower curvature bound $-A\kappa$. Its canonical limit is
quasi-plurisubharmonic. Distributional convergence gives the
current identities and strict positivity.

Fix $\pi$ and $E$, and choose a Gauduchon form $\gamma$ on the
compact surface $Z$. For the tail
$r_N=\sum_{j>N}a_j\ell_j$ of total coefficient
$A_N=\sum_{j>N}a_j$, the closed positive current
\[
 T_N=A_N\pi^*\kappa+\ddc(r_N\circ\pi)
\]
has trace mass $A_N\int_Z\pi^*\kappa\wedge\gamma$, by
distributional Stokes. The pullback of $r_N$ is
quasi-plurisubharmonic, since $\pi$ is an isomorphism on a
dense open set where $r_N$ is finite almost everywhere.
Divisorial removal gives
\[
 0\leq\nu(r_N\circ\pi,E)
 \leq A_N\frac{\int_Z\pi^*\kappa\wedge\gamma}{\int_E\gamma}.
\]
The denominator is positive and the ratio is finite.
Finite additivity of Lelong numbers now identifies the partial
profile with $\nu(u\circ\pi,E)$ up to an error tending to zero.
This proves convergence and the asserted equality.
\end{proof}

\Needspace{16\baselineskip}
\begin{proposition}\label{countable-divisorial:positive-mass}\index{Multiplier ideal!positive singularity volume}\index{Projective space!countable divisorial ideal type}
Let $v\in\PSH(X,\theta)$. The conditions
\[
 \nu(v\circ\pi,E)=b(E)\text{ for every }\pi,E,
 \qquad v\sim_{\mathcal I}u
\]
are equivalent. Under either condition,
\[
 q=P_\theta[v]_{\mathcal I}=P_\theta[u]_{\mathcal I}\sim u,
\]
and every $w\in\PSH(X,\theta)$ with $w\sim q$ satisfies
\begin{equation}\label{eq:countable-divisorial-positive-mass}
 \int_X\mathrm{MA}^{\mathrm{ext}}_\theta(w)
 =\underline{\operatorname{vol}}_\theta(q)
 =\overline{\operatorname{vol}}_\theta(q)
 =\int_XR^2>0.
\end{equation}
The background is big and $X$ has the bounded mass property.
\end{proposition}
\begin{proof}
\emph{Ideal profile and envelope.} Choose an integer $C>A$
with $C\kappa\geq\theta$. Both $u$ and $v$ are
$C\kappa$-plurisubharmonic. The projective valuative
characterization of all-weight ideal types applies in this
closed background: generic orders on all smooth projective
modifications recover point Lelong numbers by further point
blowups, and equality of those numbers is equivalent to ideal
equivalence. The stated profile therefore implies
$v\sim_{\mathcal I}u$. Conversely, ideal equivalence pulls
back under every compact modification by
\cref{lem:birational-ideal-order}; the local valuative criterion
\cref{thm:local-i-inputs} then gives equality of generic
coefficients on every such modification. The common factor
between growth against $\log|z|$ and the book's
$\log|z|^2$ convention does not affect these equalities.

The potential $u$ has complex Hessian $S-A\kappa$.
By \cref{cor:divisorial-envelope-equality}, its ordinary and
ideal envelopes agree and are ordinarily equivalent to $u$.
Ideal equivalence identifies the defining families for the
two ideal envelopes in the statement, proving the assertion
about $q$. The inequality $\theta+\ddc u\geq R>0$ proves
bigness. For any Hermitian form $h$ on $X$, choose $D$ with
$D\kappa\geq h$. Bounded $h$-plurisubharmonic potentials are
$D\kappa$-plurisubharmonic, and positivity of bounded mixed
products bounds their masses by $\int_X(D\kappa)^2=D^2$.
This proves the bounded mass property.

\emph{The bounded residual.} Fix $w\sim u$ in
$\PSH(X,\theta)$ and put $T=C\kappa+\ddc w$.
Its base divisorial coefficients are $a_j$, so divisorial
removal and passage through increasing finite sums give
\[
 K=T-S=(C-A)\kappa+\ddc h\geq0,\qquad h=w-u.
\]
The distributional difference has a canonical
quasi-plurisubharmonic representative, which is bounded by
ordinary equivalence and almost-everywhere order. Thus $K$
has bounded local potentials, and its trace charges no
pluripolar set. The current $S$ is carried by the actual pole
union $\bigcup_jH_j\subset\{w=-\infty\}$. Thus
$\langle T\rangle=K$. For every smooth real $(1,1)$-form $L$
with $\ddc L=0$, distributional Stokes gives
\[
 \int_X L\wedge\langle T\rangle
       =(C-A)\int_XL\wedge\kappa.
\]

\emph{Closed top mass and smooth correction.} Put
$T_u=C\kappa+\ddc u=S+(C-A)\kappa$.
For the finite partial sums $u_N=\sum_{j\leq N}a_j\ell_j$,
the closed non-pluripolar mass is
$(C-\sum_{j\leq N}a_j)^2$: outside the finite divisor union
the current is the smooth residual, and the measure gives
that union zero mass. Since $u_0\leq u_N$, closed-background
mass monotonicity, the closed case of
\cref{thm:volume-monotonicity}, gives
\[
 \int_X\langle T_u^2\rangle
 \leq\left(C-\sum_{j\leq N}a_j\right)^2.
\]
For the reverse inequality, locally write $S=\ddc b$ with
$b$ plurisubharmonic. The currents
$(C-A)\kappa+\ddc\max(b,-k)$ have bounded local potentials
and dominate $(C-A)\kappa$.
Their squared products, restricted to $\{b>-k\}$, therefore
dominate $(C-A)^2\kappa^2$ on that set. Letting $k$ increase
and using the local truncation definition gives
$\langle T_u^2\rangle\geq(C-A)^2\kappa^2$.
The two bounds prove equality of total masses.
Closed mass invariance under bounded changes then gives
\[
 \int_X\langle T^2\rangle=(C-A)^2.
\]

Set $L=C\kappa-\theta\geq0$, so $\ddc L=0$.
Expanding the bounded products of $\max(w,-k)$ in the two
backgrounds and restricting to $\{w>-k\}$ gives, in the limit,
\[
 \mathrm{MA}^{\mathrm{ext}}_\theta(w)
       =\langle T^2\rangle-2L\wedge\langle T\rangle+L^2.
\]
All correction measures are finite: the closed top measure
has the mass just computed, the degree-one limit is $K$, and
the smooth measure $L^2$ gives the pole set zero mass.
The same truncations are used in both backgrounds, so this
identity is an equality of Borel measures.
Integrating yields
\[
 \int_X\mathrm{MA}^{\mathrm{ext}}_\theta(w)
 =(C-A)^2-2(C-A)\int_XL\wedge\kappa+\int_XL^2
 =\int_X(\theta-A\kappa)^2.
\]
This holds throughout the ordinary type, proving both
singularity-volume identities and their strict positivity.
\end{proof}

\Needspace{16\baselineskip}
\begin{example}\label{countable-divisorial:nonclosed-example}\index{Pluriclosed metric!nonanalytic positive-volume type}
There is a nonclosed pluriclosed background on $\mathbb P^2$
with a nonanalytic ideal type satisfying
\cref{countable-divisorial:positive-mass}. Its prescribed
divisors are not contained in a single pencil, and both
singularity volumes of its ideal model are at least $1/16$.
\end{example}
\begin{proof}
In a coordinate ball choose a smooth real cutoff $\chi$ equal
to one near the origin, and extend
$\alpha=\chi|z_2|^2\mathrm dz_1$ by zero. Put
\[
 \tau=\partial\overline\alpha+\bar\partial\alpha.
\]
It is real and smooth, with $\ddc\tau=0$. Where $\chi=1$,
\[
 \mathrm d\tau=\mathrm dz_2\wedge\mathrm d\bar z_2
                  \wedge(\mathrm dz_1-\mathrm d\bar z_1)\ne0.
\]
Choose $\epsilon>0$ small enough that
$\theta=\kappa+\epsilon\tau\geq\kappa/2$.
Set $A=1/4$, $a_j=A2^{-j}$, $t_j=2^{-j}$ and
\[
 H_j=\{Z_0+t_jZ_1+t_j^2Z_2=0\}.
\]
The Vandermonde determinant shows that any three defining
linear forms are independent; hence the hyperplanes are not
in one pencil. Here $R=\theta-A\kappa\geq\kappa/4$, so
$\int_XR^2\geq1/16$. The preceding proposition gives the
mass throughout the ordinary type and both singularity volumes.
There are infinitely many positive base divisor coefficients.
An analytic singularity on a compact manifold has only
finitely many prime divisorial components, so no analytic
potential has this ideal type.
\end{proof}
